\sectioncentered*{Заключение}
\addcontentsline{toc}{section}{Заключение}

В лабораторной работе реализована универсальная рекуррентная нейронная сеть для прогнозирования числовых последовательностей, покрывающая все шесть архитектур из вариантов задания: сеть Джордана, сеть Элмана, сеть Джордана~--~Элмана, простую рекуррентную сеть (КП), сеть с долговременной кратковременной памятью (LSTM, ДКП) и сеть с управляемыми рекуррентными блоками (GRU). Скрытый слой может использовать одну из четырёх функций активации (арксинус, тангенс, ELU, Leaky ReLU), эффектор~\mdash  линейную (опционально) или ELU (для LSTM/GRU, как требует методичка). Обучение~\mdash  обратное распространение во времени с градиентным клиппингом; многошаговый прогноз~\mdash  методом скользящего окна с сохранением контекста.

Показано, что все архитектуры сходятся на тестовых рядах. Для «мягких» последовательностей (геометрическая и арифметическая прогрессии, гармонический ряд, периодическая последовательность, квадраты, нечётные) достигается практически нулевая ошибка прогноза ($\mathrm{MSEn} \lesssim 10^{-3}$), при этом простые архитектуры (Элман, КП) точнее LSTM/GRU на коротких рядах. Для быстро растущих рядов (факториал, числа Белла, $2^n$, Фибоначчи) и апериодической последовательности ошибка существенна, что объясняется ограниченностью диапазона $[-1, 1]$ и невозможностью экстраполировать сверхлинейный рост.

Проведённые эксперименты подтвердили влияние гиперпараметров: ошибка монотонно убывает с ростом размера окна; малое число контекстных нейронов ($C = 1$--$2$) предпочтительнее избыточного контекста; зануление контекста в начале эпохи и нулевая (малая) инициализация весов эффекторов устойчиво улучшают качество сетей Джордана и Джордана~--~Элмана.

\textbf{Семантическая составляющая.} Рекуррентная сеть интерпретируется как динамическая система с внутренним состоянием (контекстом), отображающая подпоследовательность в следующий отсчёт; обучение настраивает рекуррентные связи так, чтобы модель воспроизводила закономерность ряда, а многошаговый прогноз получается итеративным разворачиванием этой динамики с сохранённым состоянием.

\textbf{Прагматическая составляющая.} Получен воспроизводимый стенд на Python с единым интерфейсом всех шести архитектур, автоматической серией экспериментов (сходимость, окно, контекст, скрытый слой, эффектор, зануление контекста, начальные веса, многошаговый прогноз), сохранением метрик и графиков, пригодный для сравнения рекуррентных архитектур и подбора гиперпараметров.

\begin{thebibliography}{9}
	\addcontentsline{toc}{section}{Список использованных источников}

	\bibitem{method}
	Методические указания к лабораторной работе №5 по дисциплине <<Модели решения задач в интеллектуальных системах>>. \mdash  Минск : БГУИР.

	\bibitem{haykin}
	Хайкин С. Нейронные сети: полный курс. \mdash  2-е изд. \mdash  М. : Вильямс, 2006. \mdash  1104~с.

	\bibitem{numpy}
	NumPy documentation [Электронный ресурс]. \mdash  Режим доступа: \url{https://numpy.org/doc/}. \mdash  Дата доступа: 15.10.2026.

	\bibitem{matplotlib}
	Matplotlib documentation [Электронный ресурс]. \mdash  Режим доступа: \url{https://matplotlib.org/stable/}. \mdash  Дата доступа: 15.10.2026.
\end{thebibliography}
